\documentclass[11pt]{article}
\usepackage[margin=1in]{geometry}
\usepackage[hidelinks]{hyperref}
\usepackage{parskip}
\usepackage{enumitem}
\usepackage[sfdefault]{roboto}
\pagenumbering{gobble}

\begin{document}

{\LARGE \textbf{Case Study}}\\
\textbf{Project: LLM-Based Customer Support Chatbot (POC)}\\
\textbf{Organization:} Aptos India\\
\textbf{Role:} Machine Learning Engineer

\section*{Problem}
Support teams managing a B2B retail analytics dashboard were spending significant time handling repetitive, knowledge-heavy queries. Existing workflows had three critical issues:
\begin{itemize}[leftmargin=1.2em]
  \item High manual dependency for common support questions
  \item Inconsistent answer quality due to scattered internal documentation
  \item Limited ability to scale high-quality responses across growing customer usage
\end{itemize}

\section*{Goal}
Build an internal LLM-powered support assistant that can provide context-aware, reliable answers from internal sources while improving response quality and reducing manual effort.

\section*{System Design}
\textbf{Architecture components}
\begin{itemize}[leftmargin=1.2em]
  \item API layer for request orchestration and response delivery
  \item Retrieval pipeline over internal knowledge and analytics metadata
  \item Prompt templates with grounding constraints
  \item Evaluation and monitoring loop for output quality
  \item Containerized deployment for controlled experimentation and scale
\end{itemize}

\textbf{Core technical choices}
\begin{itemize}[leftmargin=1.2em]
  \item Python + FastAPI for service orchestration
  \item RAG workflow using vector retrieval and context selection
  \item LangChain/LangGraph-based orchestration patterns
  \item Vertex AI and managed cloud components for operational reliability
  \item Dockerized deployment integrated with MLOps practices
\end{itemize}

\section*{Execution Approach}
\textbf{1. Retrieval quality first}
I started by improving retrieval relevance so prompts received precise context. This reduced unsupported generations and improved factual grounding.

\textbf{2. Prompt versioning and evaluation}
I introduced structured prompt versions and compared them through offline tests and runtime metrics to improve consistency and tone while preserving factual correctness.

\textbf{3. Production readiness}
I containerized the pipeline and added monitoring for failure patterns, response quality drift, and operational stability to support iterative improvement.

\section*{Results}
\begin{itemize}[leftmargin=1.2em]
  \item Reduced incorrect answers by 35\% through retrieval and prompt optimization
  \item Reduced dependency on manual support teams by 30\%
  \item Enabled faster iteration and more stable deployment via containerized MLOps workflows
\end{itemize}

\section*{What This Demonstrates for AISEO}
This project demonstrates practical capability in building LLM-powered content/answer systems where quality, relevance, and operational reliability are equally important. The same approach translates well to GEO/SEO tooling: retrieval-driven answer quality, prompt-controlled output behavior, measurable evaluation loops, and scalable backend architecture.

\section*{Technology Stack}
Python, FastAPI, RAG, LangChain, LangGraph, Vertex AI, Vector DB, Docker, MLOps.

\end{document}
